\chapter{What is noise}
\label{ch:noise}

\section*{What this chapter is for}

The chapter that saves the most time, because it is about what not to study.
It has three parts: material that is real but rarely asked, material that is a
rename of something you already know, and a test for deciding about a topic
this book does not mention.

\section{Real, and rarely asked}

\begin{kittable}{@{}lXl@{}}
\textbf{Topic} & \textbf{Why it draws attention} & \textbf{Status} \\
\midrule
Fine-tuning, LoRA, PEFT & every course leads with it & niche\source{field-guide} \\
Transformer internals & it is what the field \emph{means} to outsiders & research-adjacent\source{field-guide} \\
Attention arithmetic, KV cache & teachable and feels rigorous & rarely asked\source{field-guide} \\
Quantisation, mixture-of-experts & it is in the news & rarely asked\source{field-guide} \\
Training infrastructure & it is impressive & a different job (Ch.~\ref{ch:three-jobs}) \\
\end{kittable}

The rule that covers the whole table: \reported{if the posting does not mention
fine-tuning, you are unlikely to be asked about it.}{field-guide} It
generalises to every row.

\begin{kittrap}
This material is not worthless and the trap is not that it is wrong to learn.
The trap is that it is the only part of the subject with a settled curriculum,
so it is the part that has courses \dash{} and a course is what an engineer
reaches for when starting something unfamiliar.

The cost is specific and it compounds: a month spent here is a month not spent
on retrieval failure modes, evaluation, or cost arithmetic, which are the three
surfaces most likely to decide the round. And it produces a characteristic
failure where a candidate can explain the mechanism of attention and cannot say
what to check first when the answers are subtly wrong.
\end{kittrap}

\judgement{There is one good reason to learn some of it anyway, and it is not an
interview reason: a rough model of how these systems work makes their failure
modes less surprising. A day on the mechanism is well spent. A month is the
error.}

\section{Renamed, not new}

The second category is harder to spot, because the material is genuinely useful
and the name is genuinely current. Something you already know acquires a new
label, an ecosystem of posts and courses grows around the label, and it appears
to be a subject you have not learnt.

Two examples that were traceable when this was written, and the point is the
test rather than the examples.

\textbf{Managing what goes into the context window} has a current name that
makes it sound like a discipline. The material under it \dash{} what to
retrieve, what to keep from earlier turns, how to spend a finite budget, where
in the window to put what matters \dash{} is retrieval, prompt design and
capacity planning, covered in Chapters~\ref{ch:retrieval}, \ref{ch:cost} and
\ref{ch:model}. \reported[single]{The only sources found for it as a named
interview topic were a newsletter selling a handbook on it and an
interview-preparation site.}{kalra-transition}

\textbf{Protocol standards for connecting models to tools} are real, useful,
and worth understanding if a posting names one. As an interview topic, every
source found was a content-marketing question list, with no practitioner or
primary account. Chapter~\ref{ch:agents} covers tool design without naming a
protocol, deliberately: the design questions are the same whatever the wire
format, and the wire format is an afternoon.

\section{The test}

For any topic somebody tells you is essential:

\begin{kitmethod}
\textbf{1. Who says so?} A practitioner describing their own work, a company
describing its own process, or a corpus with a stated method \dash{} against a
site whose business is preparation material. The second category recycles
itself and produces the appearance of consensus from one origin.

\textbf{2. Does it appear in postings you are actually pursuing?} The most
reliable signal available, and it is free.

\textbf{3. Is it a new idea or a new word?} Write down what you would have to
learn. If every line maps onto something you already know, it is a word.

\textbf{4. Can you trace the claim to one origin?} Several sources repeating
one leaked note, one person's account, or one vendor's post is one source. This
is the commonest failure, and Appendix~\ref{app:refused} names the instances of
it this book found.
\end{kitmethod}

\judgement{Applying this test costs about ten minutes per topic and is the
highest-return habit in this book, because the field generates new essential
topics faster than anybody can study them and most of them do not survive the
first question.}

\section{What to do with the time instead}

The ordering from Chapter~\ref{ch:gap-map}, which this chapter exists to
protect: something small built end to end and deployed; an evaluation harness
on it; retrieval and its failure modes; the cost arithmetic, out loud; and only
then anything in the table above, and only if a posting asks.

\begin{drill}
List everything you have studied or bookmarked for this transition. Mark each
against the three categories: genuinely new, renamed, or noise. Then compare
the time spent with the ordering above.

Most people find the majority of their time in the third column, and the
reason is not carelessness \dash{} it is that the third column is the part with
courses.
\end{drill}

\section*{What this chapter does not claim}

It does not claim the topics in the first table are unimportant to the field.
It claims they are rarely assessed for this role, and that the strongest
available evidence is the posting in front of you.

The judgement about renamed topics is the author's, based on failing to find
primary sources for either as a named interview subject. Absence of evidence
here is weak evidence, and if a posting names one of them, prepare for it.
